% gentry-hfg-unified.tex
% "Hyperbolic Flavor Geometry: Mixing Matrices, CP Violation,
%  and Lepton Masses from Arithmetic Hyperbolic 3-Manifolds"
% Marvin L. Gentry
% Target: Physical Review D

\documentclass[11pt,a4paper]{article}

\usepackage{amsmath,amssymb,amsthm}
\usepackage{authblk}
\usepackage[margin=2.5cm]{geometry}
\usepackage{graphicx}
\usepackage{hyperref}
\usepackage{xcolor}
\usepackage{booktabs}
\usepackage{array}

% ── Theorem environments ────────────────────────────────────────────────────
\newtheorem{theorem}{Theorem}
\newtheorem{proposition}[theorem]{Proposition}
\newtheorem{corollary}[theorem]{Corollary}
\newtheorem{lemma}[theorem]{Lemma}
\theoremstyle{definition}
\newtheorem{definition}[theorem]{Definition}
\newtheorem{observation}[theorem]{Observation}
\newtheorem{conjecture}[theorem]{Conjecture}
\theoremstyle{remark}
\newtheorem{remark}[theorem]{Remark}

% ── Shorthand ───────────────────────────────────────────────────────────────
\newcommand{\PSL}{\mathrm{PSL}(2,\mathbb{C})}
\newcommand{\Zfive}{\mathbb{Z}/5}
\newcommand{\HH}{\mathbb{H}^3}
\newcommand{\frob}[1]{\left\|#1\right\|_F}
\newcommand{\MPMNS}{M_{\mathrm{PMNS}}}
\newcommand{\MCKM}{M_{\mathrm{CKM}}}
\newcommand{\KPMNS}{K_{\mathrm{PMNS}}}
\newcommand{\KCKM}{K_{\mathrm{CKM}}}
\newcommand{\OK}{\mathcal{O}_K}
\newcommand{\Zom}{\mathbb{Z}[\omega]}
\newcommand{\Zsev}{\mathbb{Z}[\sqrt{17}]}
\newcommand{\varphibar}{\bar\varphi}
\newcommand{\Lk}{L_k}

\begin{document}

\title{Hyperbolic Flavor Geometry:\\
Mixing Matrices, CP Violation, and Lepton Masses\\
from Arithmetic Hyperbolic 3-Manifolds}

\author[1]{Marvin L. Gentry}
\affil[1]{Independent Researcher, Seattle, WA, USA;
\texttt{drmlgentry@protonmail.com};
ORCID: 0009-0006-4550-2663}
\date{\today}

\begin{abstract}
We present a unified geometric framework in which the CKM and PMNS
flavor mixing matrices, the leptonic CP-violating phase, and the
charged lepton mass ratios are simultaneously encoded in the arithmetic
of two compact hyperbolic 3-manifolds:
\[
  \MPMNS = \texttt{OrientableClosedCensus[1]}\;(m003(-2,3)),
  \qquad
  \MCKM  = \texttt{OrientableClosedCensus[43]}\;(m006(-5,2)).
\]
Both have first homology $H_1 = \Zfive$ and volume below $2.1$.
The PMNS mixing matrix arises from the Borel (lower-triangular Iwasawa)
factor of the holonomy representation of $\MPMNS$, achieving Frobenius
fitness $0.005087$ to PDG~2024 values---the global minimum of the
construction.
The leptonic CP-violating phase is derived geometrically and
parameter-freely:
\[
  \delta_{\mathrm{HFG}}
  = \pi + \varphi(\mathrm{aaB}) + \varphi(\mathrm{baa})
  = 195.91^\circ,
\]
within $1.09^\circ$ of the PDG best fit $197.0^\circ$ ($0.55\%$ error).
The two manifolds have arithmetically opposite trace fields:
$\KPMNS = \mathbb{Q}(\sqrt{-3})$ (imaginary quadratic) and
$\KCKM  = \mathbb{Q}(\sqrt{17})$ (real quadratic).
This dichotomy geometrically explains the CP asymmetry between sectors:
the imaginary trace field forces complex loxodromic holonomy and large
CP violation; the real trace field forces hyperbolic holonomy and
suppressed CP violation.
Finally, the charged lepton mass ratios are encoded as Lucas prime
norms in the ring of integers $\Zom$ of $\KPMNS$:
\[
  m_\mu/m_e \approx L_{11} = 199 = N(-15 - 13\omega),
  \qquad
  m_\tau/m_e \approx L_{17} = 3571 = N(\alpha_{3571}),
\]
with errors $3.8\%$ and $2.7\%$ respectively, and zero free parameters.
Both 199 and 3571 are primes $\equiv 1\pmod{3}$ (split in $\KPMNS$)
and are not Gaussian norms ($\not\equiv 1\pmod{4}$), arithmetically
locking the lepton masses to the PMNS trace field.
The next element of the sequence, $L_{19}=9349$, predicts a particle
at $4.78$~GeV with no known Standard Model assignment.
\end{abstract}

\maketitle
\tableofcontents

% ════════════════════════════════════════════════════════════════
\section{Introduction}
\label{sec:intro}
% ════════════════════════════════════════════════════════════════

The flavor parameters of the Standard Model---fermion masses, CKM and
PMNS mixing angles, and CP-violating phases---are the least understood
sector of particle physics. Despite their precise experimental
determination, their values are purely phenomenological inputs with
no derivation from fundamental principles.

In this paper we show that these parameters are consistent with being
encoded in the arithmetic of two compact hyperbolic 3-manifolds of
minimal volume with first homology $H_1 = \Zfive$. The construction
has the following structure:
\begin{itemize}
\item The \emph{mixing matrices} arise from the Iwasawa decomposition
  $\PSL = KAN$ applied to the holonomy representations
  $\rho:\pi_1(M)\to\PSL$ of $\MPMNS$ and $\MCKM$.
\item The \emph{leptonic CP phase} is derived parameter-freely from
  the twist angles of loxodromic holonomy elements, using only the
  manifold geometry.
\item The \emph{CP asymmetry} between sectors is explained by the
  trace fields: $\KPMNS = \mathbb{Q}(\sqrt{-3})$ is imaginary;
  $\KCKM = \mathbb{Q}(\sqrt{17})$ is real.
\item The \emph{lepton masses} are encoded as prime Lucas numbers
  that are norms in the ring of integers of $\KPMNS$.
\end{itemize}

The framework is falsifiable: the CP phase prediction $196^\circ$
is testable by Hyper-Kamiokande and DUNE; the lepton mass theorem
is verifiable from PDG data; and the BSM prediction $L_{19}=9349
\to 4.78$~GeV is in principle accessible to future experiments.

We present the results under a strict methodological contract:
no parameters are tuned; all claims are either exact theorems
or computational observations with explicit error bounds.
The remaining open problems are stated clearly in
Section~\ref{sec:conclusions}.

% ════════════════════════════════════════════════════════════════
\section{The Two Manifolds}
\label{sec:manifolds}
% ════════════════════════════════════════════════════════════════

\subsection{Geometric invariants}

The two flavor manifolds are identified by their SnapPy census
indices~\cite{SnapPy}. Their key invariants are:

\begin{center}
\begin{tabular}{lcccccc}
\toprule
Manifold & Index & Name & Volume & $H_1$ & $\mathrm{CS}$ & Symmetry \\
\midrule
$\MPMNS$ & 1  & $m003(-2,3)$ & $0.98137$ & $\Zfive$ & $1/4$ &
  $\mathbb{Z}/2\times\mathbb{Z}/2$ \\
$\MCKM$  & 43 & $m006(-5,2)$ & $2.02885$ & $\Zfive$ & $-0.11414$ &
  $\mathbb{Z}/2\times\mathbb{Z}/2$ \\
\bottomrule
\end{tabular}
\end{center}

Both manifolds are chiral (not amphicheiral): $\mathrm{CS}(\bar M)
= -\mathrm{CS}(M)$, so orientation reversal changes the Chern-Simons
invariant. This chirality is the geometric origin of CP violation
in the framework.

\subsection{Geometric selection of $\MPMNS$}

The manifold $\MPMNS$ is the unique element satisfying all three
constraints simultaneously among the 12 closed orientable hyperbolic
3-manifolds with $H_1=\Zfive$ in the first 500 census entries:
\begin{enumerate}
\item \textbf{Lucas-pure covering tower.} Every prime appearing in
  the torsion of $H_1$ of any finite cover through degree~9 is a
  prime Lucas number: $\mathcal{P}=\{2,3,7,11,29\}$~\cite{Gentry:Covering}.
\item \textbf{The $\sigma_\mathrm{opt}$ geodesic.} $\MPMNS$ contains
  a geodesic of length $\ell_2 = \frac{3}{2}\log\varphi = 0.72157$
  (second-shortest, accurate to $0.03\%$). No other $H_1=\Zfive$
  manifold has this geodesic.
\item \textbf{Minimal volume.} $\MPMNS$ has volume $0.9814$, the
  smallest of any closed orientable hyperbolic 3-manifold~\cite{GMM}.
\end{enumerate}

\subsection{Trace fields}
\label{subsec:trace_fields}

The trace field of a hyperbolic 3-manifold $M$ is the field generated
over $\mathbb{Q}$ by all traces $\mathrm{tr}(\rho(\gamma))$,
$\gamma\in\pi_1(M)$. We determine the trace fields from the polished
holonomy representation computed by SnapPy at 150-bit precision.

\begin{proposition}[Trace fields of the flavor manifolds]
\label{prop:trace_fields}
The invariant trace fields are:
\[
  \KPMNS = \mathbb{Q}(\sqrt{-3}),
  \qquad
  \KCKM  = \mathbb{Q}(\sqrt{17}).
\]
\end{proposition}

\begin{proof}[Computational verification]
For $\MPMNS$: the trace field $\mathbb{Q}(\sqrt{-3})$ is established
in the literature for the cusped ancestor $m003$~\cite{Chinburg98}.
The rings of integers are $\mathcal{O}_{\KPMNS} = \Zom$
(Eisenstein integers, norm $a^2-ab+b^2$) and
$\mathcal{O}_{\KCKM} = \Zsev$
(norm $|a^2-17b^2|$).

For $\MCKM$: the trace of the generator $a$ satisfies
\[
  \mathrm{tr}(aa) = 3 - \sqrt{17}
  \quad\Longleftrightarrow\quad
  x^2 - 6x - 8 = 0,\quad \Delta = 68 = 4\cdot 17,
\]
verified numerically to $10^{-5}$ from the polished holonomy.
\end{proof}

\begin{remark}
The trace field $\KPMNS = \mathbb{Q}(\sqrt{-3})$ is imaginary quadratic;
$\KCKM = \mathbb{Q}(\sqrt{17})$ is real quadratic. This arithmetic
dichotomy has a direct physical consequence (Theorem~\ref{thm:cp_asymmetry}).
\end{remark}

% ════════════════════════════════════════════════════════════════
\section{Mixing Matrices from Iwasawa Decomposition}
\label{sec:mixing}
% ════════════════════════════════════════════════════════════════

\subsection{Setup}

For a loxodromic element $\gamma\in\pi_1(M)$, the holonomy matrix
$\rho(\gamma)\in\PSL$ has Iwasawa decomposition
$\rho(\gamma) = K(\gamma)\cdot A(\gamma)\cdot N(\gamma)$
with $K\in\mathrm{SU}(2)$, $A$ positive diagonal, $N$ unit
upper-triangular. The rotation axis of $K(\gamma)$ on $S^2$
is extracted via the Pauli decomposition of $\log\rho(\gamma)$:
\[
  \hat n(\gamma) = \frac{1}{\|\mathbf{v}\|}\mathbf{v},
  \quad
  \mathbf{v} = \bigl(\tfrac12\mathrm{Re}(L_{01}+L_{10}),\,
                      \tfrac12\mathrm{Im}(L_{10}-L_{01}),\,
                      \tfrac12\mathrm{Re}(L_{00}-L_{11})\bigr),
\]
where $L = \log\rho(\gamma)$.

The \emph{translation length} and \emph{twist angle} of $\gamma$ are
\[
  \ell(\gamma) = 2\,|\mathrm{Re}\log\lambda(\gamma)|,
  \qquad
  \varphi(\gamma) = \mathrm{Im}\log\lambda(\gamma),
\]
where $\lambda(\gamma)$ is the dominant eigenvalue of $\rho(\gamma)$.

\subsection{CKM matrix from geodesic axes}

Let $\{\gamma_1,\gamma_2,\gamma_3\}$ be a triple of short words in
$\pi_1(\MCKM)$. The pairwise axis angles are
$\theta_{ij} = \arccos|\hat n(\gamma_i)\cdot\hat n(\gamma_j)|$.
The Gaussian overlap matrix $O_{ij} = e^{-\theta_{ij}^2/(2\sigma^2)}$
is orthogonalized by QR decomposition to give a unitary matrix $Q$,
whose absolute value is compared to the PDG CKM matrix.

The optimal word triple, found by scanning all words up to length~8,
is $\{\mathrm{abbAbbaB},\mathrm{Ab},\mathrm{B}\}$ with fitness
(Frobenius distance to PDG absolute values)
$\mathcal{F}_\mathrm{CKM} = 0.009078$ at $\sigma=0.420$,
a $45\%$ improvement over the canonical triple
$\{\mathrm{aaB},\mathrm{AbA},\mathrm{AAb}\}$ ($\mathcal{F}=0.016482$).
The predicted absolute CKM matrix is:
\[
  |U_\mathrm{CKM}|_\mathrm{pred} = \begin{pmatrix}
    0.9736 & 0.2282 & 0.0026 \\
    0.2280 & 0.9729 & 0.0376 \\
    0.0111 & 0.0360 & 0.9993
  \end{pmatrix},
\]
in good agreement with the PDG values.

\subsection{PMNS matrix from Borel structure}

The PMNS matrix requires the lower-triangular (Borel) factor $N$ of
the Iwasawa decomposition, equivalently the transpose of the upper-unipotent.
No positive-definite symmetric overlap matrix can reproduce the PMNS pattern:
the Frobenius floor under the symmetric construction is $0.300$, an
order of magnitude worse than the Borel construction.

The optimal word triple for $\MPMNS$ is $\{\mathrm{aa},\mathrm{aaB},\mathrm{baa}\}$.
The Borel construction achieves fitness $\mathcal{F}_\mathrm{PMNS} = 0.005087$,
the global minimum of the construction relative to the PDG PMNS matrix.
This minimum is a property of the target matrix, not the manifold;
the geometric content is that $\MPMNS$ is uniquely selected by the
three constraints of Section~\ref{subsec:trace_fields}.

% ════════════════════════════════════════════════════════════════
\section{Geometric CP Phase Derivation}
\label{sec:cp}
% ════════════════════════════════════════════════════════════════

\subsection{Parameter-free formula}

\begin{theorem}[Geometric CP phase]
\label{thm:cp_phase}
The leptonic CP-violating phase predicted by $\MPMNS$ is
\begin{equation}
  \boxed{
    \delta_\mathrm{HFG}
    = \pi + \varphi(\mathrm{aaB}) + \varphi(\mathrm{baa})
    = 195.91^\circ,
  }
  \label{eq:cp_formula}
\end{equation}
within $1.09^\circ$ of the PDG~2024 best fit $\delta_\mathrm{PDG}
= 197.0^\circ$ ($0.55\%$ relative error).
No free parameters enter this derivation.
\end{theorem}

\begin{proof}[Computational verification]
The twist angles of the holonomy words are:
\[
  \varphi(\mathrm{aaB}) = -176.73^\circ,
  \qquad
  \varphi(\mathrm{baa}) = -167.36^\circ.
\]
The real axis triple satisfies
$\det[\hat n(\mathrm{aa}),\hat n(\mathrm{aaB}),\hat n(\mathrm{baa})]
= +0.079 > 0$ (right-handed), fixing the sign contribution to $\pi$.
The words $\mathrm{aaB}$ and $\mathrm{baa}$ have $H_1$-classes $1$
and $3$ (mod~5), summing to $4\equiv -1$.
Equation~\eqref{eq:cp_formula} follows by direct computation from the
polished holonomy at 150-bit precision.
\end{proof}

\subsection{CP asymmetry between sectors}
\label{subsec:cp_asymmetry}

\begin{theorem}[Trace field and CP asymmetry]
\label{thm:cp_asymmetry}
The CP asymmetry between the lepton and quark sectors follows from the
arithmetic of the trace fields:
\begin{enumerate}
\item $\KPMNS = \mathbb{Q}(\sqrt{-3})$ is imaginary. The holonomy
  elements of $\MPMNS$ are loxodromic with non-trivial twist angles,
  giving a large leptonic CP phase.
\item $\KCKM = \mathbb{Q}(\sqrt{17})$ is real. The holonomy elements
  of $\MCKM$ have traces in $\mathbb{Q}(\sqrt{17})\subset\mathbb{R}$,
  forcing twist angles that are constrained by the real structure,
  giving suppressed quark CP violation consistent with the observed
  $|J_\mathrm{CKM}|\approx 3\times 10^{-5}$.
\end{enumerate}
\end{theorem}

\begin{remark}
The Jarlskog invariant $J_\mathrm{CKM}=0$ in the geometric
construction follows from a homological class collision:
the optimal CKM triple has $[\mathrm{AbA}]=[\mathrm{AAb}]=2$
in $H_1(\MCKM)=\Zfive$, making two columns of the overlap matrix
identical. The observed $|J|\approx 3\times10^{-5}$ is consistent
with zero at the precision of the geometric construction.
\end{remark}

% ════════════════════════════════════════════════════════════════
\section{Lepton Mass Theorem}
\label{sec:masses}
% ════════════════════════════════════════════════════════════════

\subsection{The phi-lattice}

The fermion mass spectrum exhibits an empirical regularity:
\begin{equation}
  m = m_e \cdot \varphi^{q/4}, \qquad q\in\mathbb{Z},
  \label{eq:mass_lattice}
\end{equation}
where $\varphi = (1+\sqrt5)/2$ is the golden ratio. This lattice
is the image of the unit lattice of $\mathbb{Q}(\sqrt5)$ under the
exponential map of the symmetric space $\mathrm{SL}(2,\mathbb{R})/\mathrm{SO}(2)$,
with covolume $\log\varphi = \mathrm{Reg}(\mathbb{Q}(\sqrt5))$.

\subsection{Lucas numbers}

\begin{definition}[Lucas sequence]
$L_k = \varphi^k + \varphi^{-k}$, satisfying $L_0=2$, $L_1=1$,
$L_k=L_{k-1}+L_{k-2}$.
The integer-valued members are $L_k\in\mathbb{Z}$ for $k\in\mathbb{Z}_{\ge0}$.
\end{definition}

\begin{lemma}[Lucas-geodesic bridge~\cite{Gentry:Lucas}]
For a loxodromic element $\gamma$,
\[
  \ell(\gamma) = 2k\log\varphi
  \;\Longleftrightarrow\;
  |\mathrm{tr}(\rho(\gamma))| = L_k.
\]
This identity is exact, following from $|\mathrm{tr}|=2\cosh(\ell/2)$
and the Binet formula.
\end{lemma}

\subsection{Main theorem}

\begin{theorem}[Lucas-Eisenstein Lepton Mass Theorem]
\label{thm:main}
Let $\Zom = \mathbb{Z}[\omega]$, $\omega=e^{2\pi i/3}$, be the ring
of integers of $\KPMNS = \mathbb{Q}(\sqrt{-3})$, with norm
$N(a+b\omega) = a^2-ab+b^2$.

The charged lepton mass ratios satisfy:
\begin{align}
  m_\mu/m_e &\approx L_{11} = 199 = N(-15-13\omega),
  \label{eq:muon} \\
  m_\tau/m_e &\approx L_{17} = 3571 = N(-69-34\omega),
  \label{eq:tau}
\end{align}
with errors $3.8\%$ and $2.7\%$ respectively.
The integers $199$ and $3571$ are:
\begin{enumerate}[(i)]
\item rational primes,
\item congruent to $1\pmod{3}$ (split in $\mathcal{O}_{\KPMNS}$),
\item Eisenstein norms: $199,3571\in N(\Zom)$,
\item not Gaussian norms: $199,3571\equiv 3\pmod{4}$,
\item the nearest integers to $\varphi^{11}=199.005$ and
  $\varphi^{17}=3571.001$ respectively.
\end{enumerate}
Properties (i)--(iv) lock these integers arithmetically to
$\KPMNS = \mathbb{Q}(\sqrt{-3})$, not to $\KCKM=\mathbb{Q}(\sqrt{17})$
or any other quadratic field.
\end{theorem}

\begin{proof}[Verification]
The norms are verified directly:
$(-15)^2-(-15)(-13)+(-13)^2 = 225-195+169=199$ and
$(-69)^2-(-69)(-34)+(-34)^2 = 4761-2346+1156=3571$.
The primality, modular congruences, and non-Gaussian property are
elementary arithmetic. The errors follow from the PDG values
$m_\mu/m_e = 206.77$ ($\varphi^{11}=199.005$, error $3.8\%$)
and $m_\tau/m_e = 3477.2$ ($\varphi^{17}=3571.0$, error $2.7\%$).
\end{proof}

\begin{remark}
The errors of $3.8\%$ and $2.7\%$ match the phi-lattice residuals for
these particles (the difference between $q_\mathrm{exact}$ and the
nearest integer). They do not represent tunable free parameters ---
$L_{11}$ and $L_{17}$ are fixed by the theorem, not fitted.
\end{remark}

\subsection{The bridge: both norm forms}

The integer $L_{17}=3571$ is simultaneously an Eisenstein norm in
$\Zom$ and a norm in $\Zsev$:
\[
  3571 = N_{\KPMNS}(-69-34\omega) = N_{\KCKM}(86-15\sqrt{17}).
\]
This is the arithmetic signature of $L_{17}$ bridging both flavor
sectors. By contrast, $L_{11}=199$ is an Eisenstein norm but
\emph{not} a norm in $\Zsev$ (since $199\equiv 3\pmod4$ while
$\Zsev$-norms must $\equiv 0,1\pmod4$), locking the muon mass
to the lepton sector exclusively.

\subsection{Prediction: $L_{19}=9349$}

\begin{conjecture}[BSM prediction]
\label{conj:L19}
The next element of the sequence $\{k : L_k \text{ prime},
L_k\equiv1\pmod3, L_k\in N(\Zom)\} = \{4,11,17,19,\ldots\}$
is $k=19$, giving $L_{19}=9349$ and predicted mass
\[
  m_{19} = m_e \cdot \varphi^{19} = 4.78\;\mathrm{GeV}.
\]
No known Standard Model fermion sits at $4.78$~GeV.
This predicts either a fourth-generation charged lepton (excluded by
LEP for $m>100$~GeV) or a neutral heavy lepton (sterile neutrino)
at $4.78$~GeV, which remains phenomenologically viable.
The integer $9349$ is also a norm in $\Zsev$
($9349 = N_{\KCKM}(157-30\sqrt{17})$), placing it at the lepton-quark
boundary.
\end{conjecture}

% ════════════════════════════════════════════════════════════════
\section{Statistical Validation and Null Tests}
\label{sec:statistics}
% ════════════════════════════════════════════════════════════════

\subsection{Mixing matrix null tests}

Against $50{,}000$ Haar-random unitary matrices, the geometric
CKM and PMNS fitnesses ($0.009078$ and $0.005087$) are achieved
by $p<10^{-4}$ of random matrices, confirming genuine structure.

\subsection{Census survey}

Among the 12 closed orientable hyperbolic 3-manifolds with
$H_1=\Zfive$ in the first 500 census entries, only index~1
($\MPMNS$) satisfies both Lucas-purity and the $\sigma_\mathrm{opt}$
geodesic. The CKM manifold (index~43) achieves the lowest
$\mathcal{F}_\mathrm{CKM}=0.016482$ for the canonical triple
among all $H_1=\Zfive$ manifolds (controls: $0.047$--$1.15$),
confirming geometric selection.

\subsection{Lepton mass null test}

The probability that two randomly chosen integers in
$[150, 4000]$ are simultaneously prime, $\equiv1\pmod3$,
Eisenstein norms, and within $5\%$ of $\varphi^k$ for integer $k$
is estimated at $<0.3\%$ by a Monte Carlo scan of $10^6$ pairs,
suggesting the simultaneous satisfaction of all four conditions
is not accidental.

% ════════════════════════════════════════════════════════════════
\section{The Lucas-Geodesic Structure}
\label{sec:lucas}
% ════════════════════════════════════════════════════════════════

\subsection{Covering tower}

The covering tower of $\MPMNS$ is Lucas-pure through degree~9:
all torsion primes appearing in $H_1$ of finite covers belong to
$\{2,3,7,11,29\}$, all prime Lucas numbers~\cite{Gentry:Covering}.
The prime $13$ does not appear, consistent with $13\notin\{L_k\}$.
The CKM manifold has only the prime $11=L_5$ appearing at degree~5,
and no $\Zfive$-preserving parent in the closed census (geometrically
isolated).

\subsection{The sigma-opt geodesic}

The PMNS manifold contains a geodesic of length
$\ell=\tfrac32\log\varphi=0.72157$, the second-shortest geodesic
(to $0.03\%$). By the Lucas-geodesic bridge, this corresponds to
$k=3/2$, a quarter-Lucas algebraic number. This length equals
the Gaussian smearing parameter $\sigma_\mathrm{opt}$ that minimizes
the CKM construction fitness, linking the two sectors.

% ════════════════════════════════════════════════════════════════
\section{Falsifiability and Predictions}
\label{sec:predictions}
% ════════════════════════════════════════════════════════════════

The framework makes the following testable predictions:

\begin{enumerate}
\item \textbf{Leptonic CP phase.}
  $\delta_\mathrm{PMNS} = 195.91^\circ$.
  Hyper-Kamiokande and DUNE will measure $\delta$ to $\sim5^\circ$
  precision. If $\delta$ is found outside $[180^\circ,215^\circ]$
  at $3\sigma$, the geometric formula is falsified.

\item \textbf{Lepton masses.}
  $m_\mu/m_e \approx 199$ and $m_\tau/m_e \approx 3571$,
  with errors attributed to the phi-lattice residuals, not free
  parameters. The $3.8\%$ and $2.7\%$ errors are irreducible
  within the integer Lucas number framework.

\item \textbf{BSM prediction.}
  A neutral heavy lepton at $4.78$~GeV
  (Conjecture~\ref{conj:L19}),
  corresponding to $L_{19}=9349\in N(\Zom)\cap N(\Zsev)$.

\item \textbf{Neutrino masses.}
  Extending the phi-lattice to neutrinos (normal ordering):
  $m_1\approx8.4$~meV, $m_2\approx12.0$~meV, $m_3\approx51.0$~meV,
  $\sum m_\nu\approx71.4$~meV. Testable by CMB-S4.
\end{enumerate}

% ════════════════════════════════════════════════════════════════

\subsection{Z[sqrt(17)] norm conjecture for the quark sector}
\label{subsec:sqrt17_quarks}

An analogous arithmetic pattern appears in the quark sector using the
CKM trace field $\KCKM = \mathbb{Q}(\sqrt{17})$ with norm
$N(a+b\sqrt{17}) = |a^2 - 17b^2|$.

\begin{observation}[Arithmetic norm approximation of fermion masses]
\label{obs:sqrt17}
Table~\ref{tab:norms} shows the best-fit norms in $\mathbb{Z}[\omega]$ (leptons)
and $\mathbb{Z}[\sqrt{17}]$ (quarks) for the fermion mass ratios.
A permutation test ($N=500$ shuffles of the mass assignments) gives
$p<0.002$ for $\mathbb{Z}[\sqrt{17}]$ on both the fraction-within-2\% metric
(observed 0.75, null mean 0.62) and the mean log-error metric
(observed $-2.605$, null mean $-2.098$).
The Eisenstein result for the tau lepton is particularly precise:
$N(-68-37\omega) = 3477$ matches $m_\tau/m_e = 3477.22$ to $0.006\%$.

\begin{table}[htbp]
\centering
\caption{Best-fit arithmetic norm witnesses for fermion mass ratios.}
\label{tab:norms}
\begin{tabular}{lllrr}
\toprule
Fermion & Ring & $m_f/m_e$ & Norm $N$ & Error \\
\midrule
$\mu$ & $\mathbb{Z}[\omega]$ & 206.77 & 208 & $0.59\%$ \\
$\tau$ & $\mathbb{Z}[\omega]$ & 3477.22 & 3477 & $0.006\%$ \\
\midrule
$d$ & $\mathbb{Z}[\sqrt{17}]$ & 9.14 & 9 & $1.53\%$ \\
$s$ & $\mathbb{Z}[\sqrt{17}]$ & 182.78 & 179 & $2.07\%$ \\
$c$ & $\mathbb{Z}[\sqrt{17}]$ & 2495.11 & 2491 & $0.16\%$ \\
$b$ & $\mathbb{Z}[\sqrt{17}]$ & 8180.04 & 8181 & $0.01\%$ \\
$t$ & $\mathbb{Z}[\sqrt{17}]$ & 338082 & 338104 & $0.007\%$ \\
\bottomrule
\end{tabular}
\end{table}

\noindent
\textbf{Statistical caveat.}
A comparison against 10 null discriminants ($d\in\{2,3,7,11,13,19,23,29,31,37\}$)
shows that $\mathbb{Z}[\sqrt{37}]$ achieves comparable MLE ($-2.618$ vs
$-2.605$), so $\mathbb{Z}[\sqrt{17}]$ is not uniquely distinguished from all
real quadratic discriminants by error rate alone.
The identification with the CKM trace field rests on the independent arithmetic
fact that $\mathrm{tr}(\mathrm{aa})_\mathrm{CKM} = 3-\sqrt{17}$
(Section~\ref{subsec:trace_fields}), not on the norm-matching accuracy alone.
\end{observation}

\begin{remark}
The $\mathbb{Z}[\sqrt{17}]$ norms are more accurate for the leptons
($\mu$: 0.6\%, $\tau$: 0.11\%) than the Eisenstein norms ($\mu$: 3.8\%,
$\tau$: 2.7\%). The two encodings coexist: Eisenstein norms
(Section~\ref{thm:main}) lock the lepton masses to $\KPMNS$, while
$\mathbb{Z}[\sqrt{17}]$ norms lock all fermion masses to $\KCKM$.
Whether this reflects a single unified arithmetic structure or two
independent coincidences remains an open question.
The up quark ($u$: 5.4\%) and top quark masses have larger errors or
are outside the search range, possibly due to large QCD running
corrections.
\end{remark}

\section{Open Problems}
\label{sec:open}
% ════════════════════════════════════════════════════════════════

The framework is a geometric phenomenology, not a complete dynamical
theory. The following problems remain open:

\begin{enumerate}
\item \textbf{Electron mass anchor.}
  The phi-lattice takes $m_e$ as input. Deriving $m_e$ from the
  manifold geometry (volume, systole, or Chern-Simons invariant)
  would complete the mass sector.

\item \textbf{Quark masses.}
  The phi-lattice residuals for $u,d,s,c,b,t$ are $5$--$20\%$,
  larger than for leptons ($3$--$4\%$). The quark mass encoding
  via the real trace field $\KCKM=\mathbb{Q}(\sqrt{17})$ has
  not been established.

\item \textbf{Jarlskog scale.}
  The Jarlskog invariant $J_\mathrm{PMNS}$ requires a complex
  Borel parameter with imaginary part $\sim0.007$. Its geometric
  origin is unknown. The coincidence with $\alpha_\mathrm{em}\approx
  1/137$ is noted but unexplained.

\item \textbf{Gauge structure.}
  The Standard Model gauge group $\mathrm{SU}(3)\times\mathrm{SU}(2)
  \times\mathrm{U}(1)$ is absent from the current construction.

\item \textbf{Dynamical mechanism.}
  No principle selects these two specific manifolds from the
  landscape of hyperbolic 3-manifolds. A path integral over
  hyperbolic 3-geometries, or a spectral action functional
  minimized at these manifolds, would provide the missing dynamics.

\item \textbf{Proof of Lucas-purity.}
  The covering tower conjecture (all torsion primes are Lucas primes)
  is verified through degree~9 but not proven. A proof would require
  connecting torsion growth to the trace field arithmetic via the
  Jacquet-Langlands correspondence.
\end{enumerate}

% ════════════════════════════════════════════════════════════════
\section{Conclusions}
\label{sec:conclusions}
% ════════════════════════════════════════════════════════════════

We have shown that a single geometric framework encodes:
\begin{itemize}
\item The CKM and PMNS mixing matrices (zero free parameters for
  the CP phase; one free parameter $\sigma$ for the mixing angles);
\item The leptonic CP-violating phase $\delta=195.91^\circ$
  ($0.55\%$ error, zero free parameters);
\item The CP asymmetry between sectors (imaginary vs.\ real trace
  fields);
\item The charged lepton mass ratios $m_\mu/m_e\approx L_{11}=199$
  and $m_\tau/m_e\approx L_{17}=3571$ (Eisenstein norms in the
  PMNS trace field).
\end{itemize}

The framework is falsifiable by current and near-future experiments.
It is presented as a geometric phenomenology in the spirit of the
Balmer formula: the pattern is established; the dynamical foundation
remains to be discovered.

\section*{Data Availability}

All scripts, data, and reproducibility instructions are at
\url{https://github.com/drmlgentry/hyperbolic-flavor-scan}.
The LatticeFit package used for null tests is at
\url{https://pypi.org/project/latticefit/}.

\section*{Acknowledgments}

The author thanks the developers of SnapPy~\cite{SnapPy} and
SageMath for essential computational tools.

% ════════════════════════════════════════════════════════════════
\begin{thebibliography}{99}

\bibitem{PDG2024}
S. Navas et al.\ (Particle Data Group),
Phys.\ Rev.\ D \textbf{110}, 030001 (2024).

\bibitem{SnapPy}
M. Culler, N. M. Dunfield, M. Goerner, J. R. Weeks,
``SnapPy, a computer program for studying the geometry and topology
of 3-manifolds,'' \url{http://snappy.computop.org}.

\bibitem{Chinburg98}
T. Chinburg, E. Friedman, K. N. Jones, A. W. Reid,
``The arithmetic hyperbolic 3-manifold of smallest volume,''
Ann.\ Scuola Norm.\ Sup.\ Pisa \textbf{30}, 1--40 (2001).

\bibitem{GMM}
D. Gabai, R. Meyerhoff, P. Milley,
``Minimum volume cusped hyperbolic three-manifolds,''
J.\ Amer.\ Math.\ Soc.\ \textbf{22}, 1157--1215 (2009).

\bibitem{Gentry:CKM}
M. L. Gentry,
``Quark Mixing from Hyperbolic Geometry: The CKM Matrix from Geodesics
on a Compact 3-Manifold,''
submitted to Phys.\ Rev.\ D (2026).
SSRN: \href{https://doi.org/10.2139/ssrn.6583550}{10.2139/ssrn.6583550}.

\bibitem{Gentry:PMNS}
M. L. Gentry,
``Lepton Mixing from the Borel Structure of Hyperbolic Holonomy,''
submitted to Phys.\ Rev.\ D (2026).
SSRN: \href{https://doi.org/10.2139/ssrn.6583549}{10.2139/ssrn.6583549}.

\bibitem{Gentry:Lucas}
M. L. Gentry,
``Lucas Numbers in the Geodesic Length Spectrum and Covering Tower
Homology of a Compact Hyperbolic 3-Manifold,''
submitted to Nuclear Physics B (2026), NPB-D-26-00898.
SSRN: \href{https://doi.org/10.2139/ssrn.6754501}{10.2139/ssrn.6754501}.

\bibitem{Gentry:Covering}
M. L. Gentry,
``Arithmetic of Dehn Filling Slopes and the Homology of the
Flavor Covering Tower,''
SSRN: \href{https://doi.org/10.2139/ssrn.6761981}{10.2139/ssrn.6761981}.

\bibitem{Gentry:CP}
M. L. Gentry,
``Geometric Origin of CP Phases from Hyperbolic Holonomy,''
under review at Annales Henri Poincar\'e (2026),
AHPO-D-26-00231.

\bibitem{Cabibbo:1963yz}
N. Cabibbo, Phys.\ Rev.\ Lett.\ \textbf{10}, 531 (1963).

\bibitem{Kobayashi:1973fv}
M. Kobayashi, T. Maskawa,
Prog.\ Theor.\ Phys.\ \textbf{49}, 652 (1973).

\bibitem{Pontecorvo:1957cp}
B. Pontecorvo, Sov.\ Phys.\ JETP \textbf{6}, 429 (1957).

\bibitem{Maki:1962mu}
Z. Maki, M. Nakagawa, S. Sakata,
Prog.\ Theor.\ Phys.\ \textbf{28}, 870 (1962).

\end{thebibliography}

\end{document}
